\section{Method}
\label{sec:method}

We study how RePaint resampling interacts with the training regime used for diffusion inpainting.
This section documents the models, training recipe, inference grid, and evaluation protocol used throughout.
Architecture and schedule details match our prior technical report~\citep{hamouda2026protocol}; we summarize the essentials here for completeness.

\subsection{Models}
\label{sec:models}
Both variants are $\varepsilon$-prediction DDPMs~\citep{ho2020ddpm} with a shared U-Net capacity recipe:
base width $64$, channel multipliers $(1,2,4)$, two residual blocks per level, self-attention at $16{\times}16$, and dropout $0.1$.
The diffusion schedule is linear with $T{=}1000$ and $\beta_t\in[10^{-4},0.02]$.
Full layer definitions and training logs appear in~\citet{hamouda2026protocol}.

\paragraph{Mask-conditioned DDPM.}
The network predicts noise from the channel-wise concatenation
$(x_t,\tilde{x}_0,m)\in\mathbb{R}^{7\times64\times64}$,
where $\tilde{x}_0{=}m\odot x_0$ is the known-pixel image and $m\in\{0,1\}^{1\times64\times64}$ marks known pixels.
This is the Palette-style conditioning path used in our prior work~\citep{hamouda2026protocol,saharia2022palette}.

\paragraph{Unconditional DDPM.}
The network receives only $x_t\in\mathbb{R}^{3\times64\times64}$ and never observes masks during training.
At test time, the mask enters solely through RePaint's noise-matched known-pixel stitch and optional resampling jumps~\citep{lugmayr2022repaint}, matching the classic RePaint setup.

\subsection{Training}
\label{sec:training}
We train on CelebA~\citep{liu2015celeba} faces: torchvision train split, center-crop $178{\times}178$, resize to $64{\times}64$, RGB in $[0,1]$.
Both models use Adam with learning rate $2{\times}10^{-4}$ and batch size $32$, seeded identically at initialization (seed $42$).\footnote{%
Training entrypoints call \texttt{torch.manual\_seed} with \texttt{seed: 42} from the released CelebA conditioned / unconditional YAML configs.}
This training seed happens to equal the evaluation base seed used later, but the two uses are independent---weight initialization and data order for training, versus eval index permutation and per-sample $x_T$ for evaluation.

The objective is standard noise-prediction MSE,
$\mathcal{L}{=}\mathbb{E}\|\hat{\varepsilon}_\theta(\cdot,t)-\varepsilon\|_2^2$,
with no hole weighting ($\lambda{=}0$).
Conditioned training draws masks online from the set \{\textrm{center}, \textrm{rectangle}, \textrm{brush}, \textrm{holes}\};
unconditional training uses clean images only.
All quantitative results in this paper use the \textbf{epoch-$30$} checkpoint of each model (not later finetunes).

\subsection{Inference: RePaint with $(j,r)$}
\label{sec:inference}
At test time we run full reverse length $T{=}1000$ (no truncated schedules).
Let $m$ mark known pixels and $x_0$ the clean image.
After a reverse step that produces a candidate hole fill $\hat{x}_{t-1}^{\mathrm{unk}}$, RePaint forms the noise-matched stitch
\begin{equation}
  x_{t-1}
  \;=\;
  m\odot q(x_0,t{-}1)
  \;+\;
  (1-m)\odot \hat{x}_{t-1}^{\mathrm{unk}},
  \label{eq:stitch}
\end{equation}
where $q(x_0,t{-}1)$ is a forward diffusion sample of the known pixels to noise level $t{-}1$~\citep{lugmayr2022repaint}.
When $r{>}1$, the sampler also applies forward jumps of length $j$: one step of the DDPM forward process
\begin{equation}
  x_{t}
  \;=\;
  \sqrt{1-\beta_{t}}\,x_{t-1}
  \;+\;
  \sqrt{\beta_{t}}\,\varepsilon,
  \qquad \varepsilon\sim\mathcal{N}(0,I),
  \label{eq:jump}
\end{equation}
then denoise and stitch again.
Setting $r{=}1$ disables these jumps and retains only~\eqref{eq:stitch}.
We do not shorten the reverse chain in the main grid, so reverse length equals the trained schedule.

\subsection{Experiment grid}
\label{sec:grid}
Unless noted, evaluation uses the CelebA validation split and center square masks whose side length is $\mathrm{cr}\cdot H$ (deterministic given $\mathrm{cr}$).
The primary grid is:
\begin{itemize}
  \item \textbf{Conditioned (epoch $30$), center masks.}
  $\mathrm{cr}\in\{0.20,0.30,0.40,0.50\}$, $r\in\{1,5,10\}$.
  Jump length is fixed at $j{=}10$ for $\mathrm{cr}\in\{0.20,0.30\}$.
  For the harder masks $\mathrm{cr}\in\{0.40,0.50\}$ we additionally sweep $j\in\{5,10,20\}$ (crossed with $r\in\{1,5,10\}$, with $r{=}1$ always evaluated at $j{=}10$).
  \item \textbf{Unconditional (epoch $30$), center masks.}
  $\mathrm{cr}\in\{0.40,0.50\}$ only (no $0.20$/$0.30$), with the same hard-mask $j{\times}r$ structure as the conditioned $\mathrm{cr}\in\{0.40,0.50\}$ rows.
  \item \textbf{Brush masks (conditioned only).}
  Brush strokes with coverage matched to a $\mathrm{cr}{=}0.40$ center hole; $j{=}10$, $r\in\{1,10\}$ (no $r{=}5$).
  \item \textbf{External qualitative check.}
  We also run RePaint's publicly released unconditional CelebA-HQ checkpoint at $256{\times}256$ with respaced $T{=}100$, $j{=}10$, $r\in\{1,5,10\}$, $n{=}8$, for qualitative comparison only (no metrics).
\end{itemize}

\subsection{Evaluation protocol and seeding}
\label{sec:eval-protocol}
Each cell uses $n{=}32$ images, batch size $1$, and base seed $42$.
Image indices are drawn once as a fixed length-$32$ permutation of the validation set under that seed, so every $(j,r,\mathrm{cr})$ cell for a given mask type shares the same identities.\footnote{%
In the released evaluation script: a seeded length-$32$ random permutation of the validation indices.}
Center masks are a deterministic function of $\mathrm{cr}$; brush masks follow the generator defaults for that run.

\paragraph{Per-sample reseeding.}
All paired ablations reseed the RNG before each image.
Before inpainting sample $i\in\{0,\ldots,31\}$, we set the RNG seed to $42{+}i$.\footnote{%
Implemented as a per-sample reseeding option in the released evaluation script (batch size $1$), so paired $r$/$j$ runs share the same $x_T$.}
This makes the initial $x_T$ (and the stochastic reverse/jump noise stream keyed off that seed) identical across runs that differ only in $r$ or $j$, which is required for paired $\Delta$ statistics.
Without reseeding, a longer $r$ would advance the global RNG and desynchronize later images.

\subsection{Metrics}
\label{sec:metrics}
For every image we record image-wide PSNR, hole-only PSNR, SSIM~\citep{wang2004ssim}, and LPIPS~\citep{zhang2018lpips} (AlexNet).
Headline significance tests in Section~\ref{sec:experiments} (Table~\ref{tab:resampling-delta}) use paired $\Delta$PSNR and $\Delta$LPIPS.
We omit hole-PSNR from that table because, under our protocol, only the masked region changes between paired conditions, so hole-PSNR and image-PSNR differences are numerically equivalent in the paired tests.
SSIM is tracked throughout and retained in the released per-image summaries, but is not used for the main significance analysis.
